% Part: first-order-logic
% Chapter: beyond
% Section: intuitionistic-logic

\documentclass[../../../include/open-logic-section]{subfiles}

\begin{document}

\olfileid{fol}{byd}{il}

\olsection{అంతఃప్రజ్ఞావాద తర్కం}

ద్వితీయ-స్థాయి, ఉన్నత-స్థాయి తర్కాలకు భిన్నంగా, అంతఃప్రజ్ఞావాద
మొదటిస్థాయి తర్కం సాంప్రదాయిక రూపానికి ఒక పరిమితి; మరింత
``నిర్మాణాత్మకమైన'' హేతుచింతనను నమూనీకరించడానికి దీనిని
ఉద్దేశించారు. దీని వెనుక ఉన్న కొన్ని ప్రేరణలను కింది ఉదాహరణలు
వివరిస్తాయి.

ఒక రోజు ఎవరైనా మీ దగ్గరకు వచ్చి, కింది ధర్మం గల సహజ సంఖ్య~$x$ను
నిర్ణయించానని ప్రకటించారనుకోండి: $x$ ప్రధాన సంఖ్య అయితే రీమాన్
పరికల్పన సత్యం; $x$ సంయుక్త సంఖ్య అయితే రీమాన్ పరికల్పన అసత్యం.
అద్భుతమైన వార్త!{} రీమాన్ పరికల్పన సత్యమా కాదా అనేది గణితంలో పెద్ద
అపరిష్కృత ప్రశ్నల్లో ఒకటి; ఇక్కడ వారు సమస్యను ఒక గణనకు, అంటే ఒక
నిర్దిష్ట సంఖ్య ప్రధాన సంఖ్యా కాదా అని నిర్ణయించడానికి తగ్గించినట్లు
కనిపిస్తోంది.

$x$ యొక్క ఆ అద్భుత విలువ ఏమిటి? వారు దాన్ని ఇలా వివరిస్తారు: రీమాన్
పరికల్పన సత్యమైతే $7$కు, కాకపోతే $9$కు సమానమైన సహజ సంఖ్యే~$x$.

కోపంతో మీరు చెల్లించిన డబ్బు తిరిగి ఇవ్వాలని అడుగుతారు. సాంప్రదాయిక
దృక్కోణం నుంచి పైనున్న వివరణ నిజంగానే $x$కు ఒకే విలువను
నిర్ణయిస్తుంది; కానీ మీకు నిజంగా కావలసింది \emph{స్పష్టంగా} ఇచ్చిన
$x$ విలువ.

మరొక, బహుశా తక్కువ కల్పితమైన ఉదాహరణగా కింది ప్రశ్నను పరిశీలించండి.
ఒక అకరణీయ సంఖ్యను కరణీయ ఘాతానికి హెచ్చించి కరణీయ ఫలితాన్ని పొందడం
సాధ్యమని మనకు తెలుసు. ఉదాహరణకు, $\sqrt{2}^2 = 2$. ఒక అకరణీయ సంఖ్యను
\emph{అకరణీయ} ఘాతానికి హెచ్చించి కరణీయ ఫలితాన్ని పొందడం సాధ్యమా
అనేది అంత స్పష్టం కాదు. కింది సిద్ధాంతం దీనికి అవునని సమాధానం
ఇస్తుంది:

\begin{thm}
$a^b$ కరణీయమయ్యే అకరణీయ సంఖ్యలు $a$, $b$ ఉన్నాయి.
\end{thm}

\begin{proof}
$\sqrt{2}^{\sqrt{2}}$ను పరిశీలించండి. ఇది కరణీయమైతే మన పని
పూర్తయింది: $a = b = \sqrt{2}$గా తీసుకోవచ్చు. కాకపోతే అది
అకరణీయం. అప్పుడు
\[
(\sqrt{2}^{\sqrt{2}})^{\sqrt{2}} = \sqrt{2}^{\sqrt{2} \cdot
  \sqrt{2}} = \sqrt{2}^2 = 2,
\]
ఇది కచ్చితంగా కరణీయమే. కాబట్టి ఈ సందర్భంలో $a$ను
$\sqrt{2}^{\sqrt{2}}$గా, $b$ను~$\sqrt 2$గా తీసుకోండి.
\end{proof}

ఇది చెల్లుబాటైన నిరూపణేనా? అవునని చాలామంది గణితశాస్త్రవేత్తలు
భావిస్తారు. అయినా ఇక్కడ మళ్లీ కొంత అసంతృప్తి కలిగించే విషయం ఉంది:
ఒక నిర్దిష్ట ధర్మం గల వాస్తవ సంఖ్యల జత ఉందని నిరూపించాం; కానీ అది
\emph{ఏ} జతో చెప్పలేకపోయాం. అదే ఫలితాన్ని, $a$, $b$ జతను నిరూపణలోనే
\emph{ఇచ్చే} విధంగా నిరూపించవచ్చు: $a = \sqrt{3}$,
$b = \log_3 4$గా తీసుకోండి. అప్పుడు
\[
a^b = \sqrt{3}^{\log_3 4} = 3^{1/2 \cdot \log_3 4} = (3^{\log_3
  4})^{1/2} = 4^{1/2}= 2,
\]
ఎందుకంటే $3^{\log_3 x} = x$.

మొదటి నిరూపణలోని అడుగుల వంటివాటిని అనుమతించని హేతుచింతనను
నమూనీకరించడానికి అంతఃప్రజ్ఞావాద తర్కాన్ని రూపొందించారు. $!A(x)$ను
సంతృప్తిపరచే $x$ ఉందని నిరూపించడం అంటే, రెండవ నిరూపణలోలాగా ఒక
నిర్దిష్ట~$x$ను, అది $!A$ను సంతృప్తిపరుస్తుందనే నిరూపణను ఇవ్వాలి.
$!A$ లేదా $!B$ సత్యమని నిరూపించాలంటే, వాటిలో ఒకదాన్ని
నిరూపించగలగాలి.

ఔపచారికంగా చెప్పాలంటే, మొదటిస్థాయి తర్కపు \tetoken{ఒక వ్యుత్పత్తి}{!!a{derivation}} వ్యవస్థను
ఒక నిర్దిష్ట విధంగా పరిమితం చేస్తే వచ్చేదే అంతఃప్రజ్ఞావాద
మొదటిస్థాయి తర్కం. అలాగే ద్వితీయ-స్థాయి, ఉన్నత-స్థాయి తర్కాలకు
అంతఃప్రజ్ఞావాద రూపాలు ఉన్నాయి. గణిత దృక్కోణం నుంచి ఇవి కేవలం
ఔపచారిక నిగమన వ్యవస్థలు; కానీ ఇప్పటికే గమనించినట్లు, ఒక రకమైన గణిత
హేతుచింతనను నమూనీకరించడానికి వీటిని ఉద్దేశించారు. గణితం గురించి ఒక
నిర్దిష్ట తాత్త్విక దృక్కోణం (బ్రౌవర్ అంతఃప్రజ్ఞావాదం వంటిది)
సమర్థించే హేతుచింతన ఇదని తీసుకోవచ్చు; మరింత ``మూర్తమైన'', సంతృప్తికరమైన
గణిత హేతుచింతన (బిషప్ నిర్మాణవాద పద్ధతిలోనిది) అని తీసుకోవచ్చు;
ఔపచారిక వివరణ అనౌపచారిక ప్రేరణను నిజంగా పట్టుకుంటుందా అని వాదించవచ్చు.
కానీ మన తాత్త్విక వైఖరులు ఏవైనా, అంతఃప్రజ్ఞావాద తర్కాన్ని ఔపచారికంగా
సమర్పించిన తర్కంగా అధ్యయనం చేయవచ్చు; ఏ కారణాల వల్లనైనా, అనేకమంది
గణిత తర్కశాస్త్రవేత్తలు అలా చేయడాన్ని ఆసక్తికరంగా భావిస్తారు.

అంతఃప్రజ్ఞావాద సంయోజకాలకు ఒక అనౌపచారిక నిర్మాణాత్మక అర్థనిర్దేశం
ఉంది. బ్రౌవర్, హైటింగ్, కొల్మొగొరోవ్‌ల పేర్లతో దాన్ని సాధారణంగా BHK
అర్థనిర్దేశం అంటారు. అది ఇలా ఉంటుంది: $!A \land !B$ నిరూపణలో $!A$
నిరూపణతో జతచేసిన $!B$ నిరూపణ ఉంటుంది; $!A \lor !B$ నిరూపణలో $!A$
నిరూపణ లేదా $!B$ నిరూపణ ఉంటుంది, ఏది ఉన్నదో స్పష్టమైన సమాచారం కూడా
ఉంటుంది; $!A \lif !B$ నిరూపణలో $!A$ నిరూపణను~$!B$ నిరూపణగా మార్చే
ప్రక్రియ ఉంటుంది; $\lforall[x][!A(x)]$ నిరూపణలో $x$ యొక్క ఏ
విలువకైనా $!A(x)$ నిరూపణను తిరిగి ఇచ్చే ప్రక్రియ ఉంటుంది;
$\lexists[x][!A(x)]$ నిరూపణలో $x$ యొక్క ఒక విలువతోపాటు, ఆ విలువ
$!A$ను సంతృప్తిపరుస్తుందనే నిరూపణ ఉంటుంది. ఈ అర్థనిర్దేశాన్ని
``కరీ--హోవర్డ్ సమరూపత'' లేదా ``\tetoken{సూత్రాలు}{!!{formula}s}-టైపుల నమూనా''గా పిలిచే
గణనాత్మక పదాల్లో వివరించవచ్చు: \tetoken{ఒక సూత్రం}{!!a{formula}} ఒక నిర్దిష్ట రకమైన డేటా
టైపును నిర్దేశిస్తుందని, నిరూపణలు ఆ డేటా టైపుల గణనాత్మక వస్తువులని
భావించండి; వాటివల్ల అనుగుణమైన \tetoken{సూత్రం}{!!{formula}} సత్యమని చూడగలం.

అంతఃప్రజ్ఞావాద తర్కాన్ని తరచుగా, బహిష్కృత మధ్యమ సూత్రాన్ని
``తీసివేసిన'' సాంప్రదాయిక తర్కంగా భావిస్తారు. కింది సిద్ధాంతం దీన్ని
మరింత ఖచ్చితంగా చెబుతుంది.
\begin{thm}
అంతఃప్రజ్ఞావాద తర్కంలో కింది స్వీకృత పథకాలు తుల్యమైనవి:
\begin{enumerate}
\item $(\lnot !A \lif \lfalse) \lif !A$.
\item $!A \lor \lnot !A$
\item $\lnot \lnot !A \lif !A$
\end{enumerate}
\end{thm}
మిగతా రెండింటిలో దేనినుంచైనా ఒక పథకపు నిదర్శనాలను పొందడం,
అంతఃప్రజ్ఞావాద తర్కంలో మంచి అభ్యాసం.

బ్రౌవర్ అంతఃప్రజ్ఞావాదానికి ఔపచారిక రూపాలుగా ప్రతిపాదించిన తొలి
అంతఃప్రజ్ఞావాద ప్రతిజ్ఞావాక్య తర్క నిగమన వ్యవస్థలను కొల్మొగొరోవ్,
గ్లివెంకో, హైటింగ్‌లు పరస్పరం స్వతంత్రంగా ఇచ్చారు. అంతఃప్రజ్ఞావాద
మొదటిస్థాయి తర్కానికి (అలాగే అంతఃప్రజ్ఞావాద గణితంలోని కొన్ని
భాగాలకు) తొలి ఔపచారిక రూపాన్ని హైటింగ్ ఇచ్చాడు. సాంప్రదాయికంగా
చెల్లుబాటయ్యే అనేక పథకాలు అంతఃప్రజ్ఞావాద తర్కంలో చెల్లుబాటు కాకపోయినా,
చాలావరకు చెల్లుబాటవుతాయి.

\emph{ద్వినిషేధ అనువాదం} సాంప్రదాయిక తర్కానికీ అంతఃప్రజ్ఞావాద
తర్కానికీ మధ్య ముఖ్యమైన సంబంధాన్ని వివరిస్తుంది. దాన్ని కింది విధంగా
ఆగమన పద్ధతిలో నిర్వచిస్తాం (సాంప్రదాయిక \tetoken{సూత్రం}{!!{formula}}~$!A$కు
``అంతఃప్రజ్ఞావాద'' అనువాదంగా $!A^N$ను భావించండి):
\begin{align*}
!A^N & \ident \lnot\lnot !A \quad \text{పరమాణు \tetoken{సూత్రాలు}{!!{formula}s} $!A$కు} \\
(!A \land !B)^N & \ident (!A^N \land !B^N) \\
(!A \lor !B)^N & \ident  \lnot\lnot (!A^N \lor !B^N) \\
(!A \lif !B)^N & \ident (!A^N \lif !B^N) \\
(\lforall[x][!A])^N & \ident \lforall[x][!A^N] \\
(\lexists[x][!A])^N & \ident \lnot\lnot\lexists[x][!A^N]
\end{align*}
ప్రతిజ్ఞావాక్య తర్కానికి ఈ అనువాద రూపాలను కొల్మొగొరోవ్, గ్లివెంకోలు
ఇచ్చారు; విధేయ తర్కానికి గోడెల్, గెంట్జెన్‌లు పరస్పరం స్వతంత్రంగా
ఇచ్చారు. మనకు కింది సిద్ధాంతం ఉంది.

\begin{thm}
\begin{enumerate}
\item $!A \liff !A^N$ సాంప్రదాయికంగా నిరూపించదగినది.
\item $!A$ సాంప్రదాయికంగా నిరూపించదగినదైతే, $!A^N$
  అంతఃప్రజ్ఞావాద తర్కంలో నిరూపించదగినది.
\end{enumerate}
\end{thm}

ఇప్పుడు కింది సంభాషణను ఊహించవచ్చు. సాంప్రదాయిక గణితశాస్త్రవేత్త:
``నేను $!A$ను నిరూపించాను!'' అంతఃప్రజ్ఞావాద గణితశాస్త్రవేత్త: ``మీ
నిరూపణ చెల్లదు. మీరు నిజంగా నిరూపించింది $!A^N$.'' సాంప్రదాయిక
గణితశాస్త్రవేత్త: ``నాకైతే అదే చాలు!'' ఆ రెండూ తుల్యమైనవి కాబట్టి,
సాంప్రదాయిక గణితశాస్త్రవేత్త దృష్టిలో అంతఃప్రజ్ఞావాది అనవసరంగా సూక్ష్మ
భేదం చూపుతున్నాడు. కానీ తేడా ఉందని అంతఃప్రజ్ఞావాది పట్టుబడతాడు.

పైనున్న అనువాదం శుద్ధ తర్కానికి మాత్రమే సంబంధించినదని గమనించండి;
సాంప్రదాయిక, అంతఃప్రజ్ఞావాద గణితాలకు తగిన \emph{అతార్కిక}
స్వీకృతాలు ఏవి, వాటి మధ్య సంబంధం ఏమిటి అనే ప్రశ్నలను అది
పరిశీలించదు. అయితే పైనున్న సిద్ధాంతానికి కింది స్వల్ప విస్తరణ కొంత
ఉపయోగకరమైన సమాచారాన్ని ఇస్తుంది:

\begin{thm}
$\Gamma$ సాంప్రదాయికంగా $!A$ను నిరూపిస్తే, $\Gamma^N$
అంతఃప్రజ్ఞావాద తర్కంలో $!A^N$ను నిరూపిస్తుంది.
\end{thm}

మరో మాటలో, కొన్ని పరికల్పనల నుంచి $!A$ సాంప్రదాయికంగా
నిరూపించదగినదైతే, వాటి ద్వినిషేధ అనువాదాల నుంచి $!A^N$
నిరూపించదగినది.

ఒక వాక్యం లేదా ప్రతిజ్ఞావాక్య \tetoken{సూత్రం}{!!{formula}} అంతఃప్రజ్ఞావాద తర్కంలో
చెల్లుబాటవుతుందని చూపడానికి ఒక నిరూపణ ఇస్తే చాలు. కానీ అది చెల్లదని
ఎలా చూపగలం? అందుకు యథార్థమైన, వీలైతే సంపూర్ణమైన అర్థవిచారం కావాలి.
క్రిప్కె ఇచ్చిన అర్థవిచారం ఈ అవసరాన్ని చక్కగా తీరుస్తుంది.

సాంప్రదాయిక తర్కంలో చేసిన పనినే మళ్లీ చేయవచ్చు: అర్థవిచారాన్ని
నిర్వచించి, యథార్థతను, సంపూర్ణతను నిరూపించాలి. అయితే కింది తేడాను
గమనించడం ఉపయోగకరం. సాంప్రదాయిక తర్కంలో అర్థవిచారం ఒక అర్థంలో
సంయోజకాల అర్థంలోనే అంతర్లీనంగా ఉన్న ``స్పష్టమైన'' అర్థవిచారం.
క్రిప్కె అర్థవిచారానికి కొంత అంతఃప్రజ్ఞాత్మక ప్రేరణ ఇవ్వగలిగినా, అది
అదే అనివార్యత భావాన్ని ఇవ్వదు. అంతేకాక, సాంప్రదాయిక \tetoken{నిర్మాణం}{!!{structure}}
భావన సహజమైన గణిత భావన. అందువల్ల \tetoken{ఒక నిర్మాణం}{!!a{structure}} భావనను సాంప్రదాయిక
మొదటిస్థాయి తర్కాన్ని అధ్యయనం చేసే సాధనంగా తీసుకోవచ్చు; లేదా
సాంప్రదాయిక మొదటిస్థాయి తర్కాన్ని గణిత \tetoken{నిర్మాణాలను}{!!{structure}s} అధ్యయనం చేసే
సాధనంగా తీసుకోవచ్చు. దీనికి భిన్నంగా, క్రిప్కె \tetoken{నిర్మాణాలను}{!!{structure}s}
తార్కిక నిర్మితులుగానే చూడగలం; వాటికి స్వతంత్ర గణిత ఆసక్తి ఉన్నట్లు
కనిపించదు.

ప్రతిజ్ఞావాక్య భాషకు క్రిప్కె \tetoken{నిర్మాణం}{!!{structure}}
$\mModel M = \tuple{W, R, V}$లో ఒక సమితి~$W$, ఒక కనిష్ఠ
\tetoken{మూలకం}{!!{element}} గల $W$పై పాక్షిక క్రమం~$R$, $W$ యొక్క \tetoken{మూలకాలకు}{!!{element}s}
ప్రతిజ్ఞావాక్య చరాలను ``క్రమ-ఏకదిశగా (మోనోటోన్‌గా)'' నిర్దేశించే
కేటాయింపు ఉంటాయి.
అంతఃప్రజ్ఞ ఏమిటంటే, $W$ యొక్క \tetoken{మూలకాలు}{!!{element}s} ``లోకాలను'' లేదా
``జ్ఞానస్థితులను'' సూచిస్తాయి; $v \geq u$ అనే మూలకం $u$కు
``సాధ్యమైన భవిష్యత్ స్థితి''ని సూచిస్తుంది; $u$కు నిర్దేశించిన
ప్రతిజ్ఞావాక్య చరాలు, స్థితి~$u$లో సత్యమని తెలిసిన
ప్రతిజ్ఞావాక్యాలు. ఫోర్సింగ్ సంబంధం $\mSat{M}{!A}[w]$ ఈ సంబంధాన్ని
భాషలోని యథేచ్ఛ \tetoken{సూత్రాలకు}{!!{formula}s} విస్తరిస్తుంది; $\mSat{M}{!A}[w]$ను
``స్థితి~$w$లో $!A$ సత్యం'' అని చదవండి. ఈ సంబంధాన్ని కింది విధంగా
ఆగమన పద్ధతిలో నిర్వచిస్తాం:
\begin{enumerate}
\item $\mSat{M}{\Obj p_i}[w]$ అవుతుంది అప్పుడే, $\Obj p_i$ అనేది
  $w$కు నిర్దేశించిన ప్రతిజ్ఞావాక్య చరాల్లో ఒకటైనప్పుడు.
\item $\mSat/{M}{\lfalse}[w]$.
\item $\mSat{M}{(!A \land !B)}[w]$ అవుతుంది అప్పుడే,
  $\mSat{M}{!A}[w]$, $\mSat{M}{!B}[w]$ రెండూ అయినప్పుడు.
\item $\mSat{M}{(!A \lor !B)}[w]$ అవుతుంది అప్పుడే,
  $\mSat{M}{!A}[w]$ లేదా $\mSat{M}{!B}[w]$ అయినప్పుడు.
\item $\mSat{M}{(!A \lif !B)}[w]$ అవుతుంది అప్పుడే, $w' \geq w$,
  $\mSat{M}{!A}[w']$ అయినప్పుడల్లా $\mSat{M}{!B}[w']$ అయినప్పుడు.
\end{enumerate}
ఒక ప్రతినిదర్శనాన్ని ఇచ్చే క్రిప్కె \tetoken{నిర్మాణాన్ని}{!!{structure}} రూపొందించి,
$\lnot (p \land q) \lif (\lnot p \lor \lnot q)$ అంతఃప్రజ్ఞావాద
తర్కంలో చెల్లదని చూపడానికి ప్రయత్నించడం మంచి అభ్యాసం.

\end{document}
